\section{Killed-capacity packing for hidden corridors}
\label{sec:capacity-packing}

The preceding mass bound charges vertices that are significant in one fixed
instance.  A lower-bound distribution may instead contain many decoy
corridors and make a different one relevant in each counterfactual graph.
The following capacity argument still limits every construction in which
those corridors coexist as internally disjoint passive paths.

\begin{definition}[Killed seed capacity]
\label{def:killed-capacity}
For a single seed $v$, define
\begin{equation}
 \operatorname{cap}_\alpha(v)
 :=\min_{f\in\mathbb R^V:\,f_v=1} f^\top H_\alpha f.
 \label{eq:killed-capacity}
\end{equation}
\end{definition}

The quadratic has the electrical form
\begin{equation}
 f^\top H_\alpha f
 =\alpha\sum_i d_i f_i^2
  +\frac{1-\alpha}{2}
    \sum_{\{i,j\}\in E}(f_i-f_j)^2.
 \label{eq:killed-energy}
\end{equation}
The first term is distributed killing and the second is ordinary edge
energy.

\begin{lemma}[Capacity--PPR identity; \ProvedStatus]
\label{lem:capacity-ppr}
For exact single-seed PPR,
\begin{equation}
 y_v^0=\frac{\alpha}{\operatorname{cap}_\alpha(v)}.
 \label{eq:capacity-ppr}
\end{equation}
\end{lemma}

\begin{proof}
For every symmetric positive-definite matrix $H$, minimizing
$f^\top Hf$ subject to $f_v=1$ gives value
$1/(e_v^\top H^{-1}e_v)$.  Apply this to $H_\alpha$ and use
$y^0=\alpha H_\alpha^{-1}e_v$ from \cref{eq:killed-system}.
\end{proof}

\begin{theorem}[Disjoint-corridor capacity packing; \ProvedStatus]
\label{thm:corridor-capacity}
Assume $0<\alpha\leq1/4$.  Suppose $G$ contains paths
\begin{equation}
 P_j=(v=v_{j,0},v_{j,1},\ldots,v_{j,L_j}),
 \qquad 1\leq j\leq M,
 \label{eq:corridor-paths}
\end{equation}
that are vertex-disjoint outside the seed.  Then
\begin{equation}
 \operatorname{cap}_\alpha(v)
 \geq
 \frac1{32}\sum_{j=1}^M
 \min\{\alpha L_j,\sqrt\alpha\}.
 \label{eq:corridor-capacity-bound}
\end{equation}
Consequently, if any exact normalized PPR coordinate exceeds $\tau$, then
\begin{equation}
 \sum_{j=1}^M\min\{\alpha L_j,\sqrt\alpha\}
 <\frac{32\alpha}{\tau}.
 \label{eq:corridor-packing}
\end{equation}
In particular, if every $L_j\leq1/\sqrt\alpha$, then
\begin{equation}
 \sum_{j=1}^M L_j<\frac{32}{\tau}.
 \label{eq:short-corridor-packing}
\end{equation}
\end{theorem}

\begin{proof}
Fix any $f$ with $f_v=1$ and put
\[
 K_j:=\min\left\{L_j,
                 \left\lfloor\frac1{\sqrt\alpha}\right\rfloor\right\}.
\]
Assign to path $j$ one unit of the seed's degree-weighted killing term, the
killing terms of its first $K_j$ private vertices, and the first $K_j$ edge
energies.  These assignments do not overlap: the paths are disjoint outside
$v$, their first edges are distinct, and $d_v\geq M$.  Hence
\begin{equation}
 f^\top H_\alpha f\geq\sum_{j=1}^M E_j(f),
 \label{eq:path-energy-sum}
\end{equation}
where, writing $a_k=f_{v_{j,k}}$,
\[
 E_j(f):=
 \alpha\sum_{k=0}^{K_j}a_k^2
 +\frac{1-\alpha}{2}
  \sum_{k=0}^{K_j-1}(a_k-a_{k+1})^2.
\]
If $|a_k|\geq1/2$ for every $k\leq K_j$, the killing term is at least
$\alpha K_j/4$.  Otherwise choose $m\leq K_j$ with $|a_m|<1/2$.
Since $a_0=1$, Cauchy--Schwarz gives
\[
 \sum_{k=0}^{m-1}(a_k-a_{k+1})^2
 \geq\frac{(a_0-a_m)^2}{m}
 >\frac1{4K_j}.
\]
Because $\alpha K_j^2\leq1$ and $(1-\alpha)/2\geq3/8$, both cases imply
$E_j(f)\geq\alpha K_j/16$.  Finally,
$K_j=L_j$ when $L_j\leq\lfloor1/\sqrt\alpha\rfloor$, while otherwise
$K_j\geq1/(2\sqrt\alpha)$; therefore
\[
 E_j(f)\geq\frac1{32}
             \min\{\alpha L_j,\sqrt\alpha\}.
\]
Minimizing over $f$ proves \cref{eq:corridor-capacity-bound}.

If some $y_i^0>\tau$, \cref{lem:strict-ascent} gives $y_v^0>\tau$.
Then \cref{lem:capacity-ppr} gives
$\operatorname{cap}_\alpha(v)<\alpha/\tau$, which yields
\cref{eq:corridor-packing}.  The final claim follows by substituting
$\min\{\alpha L_j,\sqrt\alpha\}=\alpha L_j$.
\end{proof}

For equal diffusion-depth corridors $L_j=\Theta(1/\sqrt\alpha)$,
\cref{eq:corridor-packing} says
\begin{equation}
 M=O(\sqrt\alpha/\tau),
 \qquad
 ML=O(1/\tau).
 \label{eq:diffusion-cancellation}
\end{equation}
This is the exact cancellation behind the path bundle in the companion
complete hybrid-local-solver note.  Its parameters are
$M=\Theta(\sqrt\alpha/\rho)$ and
$L=\Theta(1/\sqrt\alpha)$.  Moving a frontier across every arm once costs
$ML=\Theta(1/\rho)$; the proved
$\Omega(1/(\rho\sqrt\alpha))$ work comes from charging the accumulated
prefix for $L$ rounds, namely $ML^2$, under persistent support.

\begin{corollary}[No-go for disjoint hidden-port constructions;
\ProvedStatus]
\label{cor:hidden-ports}
Consider a family of nonnegative PPR instances that all contain the same
internally disjoint seed corridors with $L_j\leq1/\sqrt\alpha$, while their
completions beyond the corridor endpoints may differ.  If any member has an
exact normalized coordinate larger than $\tau$, then the total one-pass
corridor exposure is less than $32/\tau$.  In particular, changing a passive
terminal gadget, trap, or degree-preserving edge swap beyond one of these
ports cannot turn this template into an
$\Omega(1/(\tau\sqrt\alpha))$ information lower bound.
\end{corollary}

The corollary does not settle overlapping corridors, global encodings, or the
cost of rejecting counterfactual completions.  It does prove that the most
natural direct-sum construction fails for structural, rather than merely
constant-factor, reasons.
